\documentclass[12pt,a4paper]{article}
\usepackage{amsmath,amssymb,amsthm}
\usepackage{hyperref}
\usepackage{geometry}
\geometry{margin=2.5cm}
\usepackage{enumitem}
\usepackage{booktabs}

\newtheorem{theorem}{Theorem}[section]
\newtheorem{lemma}[theorem]{Lemma}
\newtheorem{corollary}[theorem]{Corollary}
\newtheorem{proposition}[theorem]{Proposition}
\theoremstyle{definition}
\newtheorem{definition}[theorem]{Definition}
\newtheorem{axiom_rs}[theorem]{Axiom}
\theoremstyle{remark}
\newtheorem{remark}[theorem]{Remark}

\title{Determinism and Apparent Randomness Reconciled\\
from Cost Minimization}

\author{Jonathan Washburn\\
Recognition Science Research Institute\\
Austin, Texas\\
\texttt{washburn.jonathan@gmail.com}}

\date{March 2026}

\begin{document}
\maketitle

\begin{abstract}
The determinism debate---whether the universe is fundamentally deterministic
(Laplace) or irreducibly random (Copenhagen)---has persisted for over a
century.  We show that the cost functional
$J(x) = \tfrac{1}{2}(x + x^{-1}) - 1$ resolves the debate by proving
both sides correct about different things.  The strict convexity of~$J$
on $\mathbb{R}_{>0}$ guarantees a unique minimizer for every constrained
optimization, making ledger dynamics deterministic at the ontological level.
Simultaneously, any finite-resolution observer accesses the ledger only
through a lossy projection map whose fibers contain multiple distinct
states, making observations \emph{appear} random.  We prove that:
(i)~the next ledger state is uniquely forced at each tick;
(ii)~the projection map is provably many-to-one;
(iii)~the fiber structure of the projection reproduces the Born-rule
probability distribution without postulating it;
(iv)~Bell's theorem is compatible because the ledger is global, not local;
and (v)~genuine decision points exist at the Gap-45 barrier, providing a
compatibilist account of free will.
\end{abstract}

%% ============================================================
\section{Introduction}
%% ============================================================

Is the universe deterministic?  The question has divided physics since
Laplace's famous claim~\cite{Laplace1814} that a sufficiently powerful
intellect could predict the entire future from the present state.
Quantum mechanics, as interpreted by the Copenhagen school, answered
no: measurement outcomes are irreducibly probabilistic
(Born rule~\cite{Born1926}), and Bell's theorem rules out local
hidden-variable completions~\cite{Bell1964}.

Yet the debate has never been fully settled.  Bohmian mechanics restores
determinism at the cost of non-locality~\cite{Bohm1952}.  Many-worlds
interpretations preserve determinism at the cost of ontological
proliferation~\cite{Everett1957}.  Collapse theories introduce
irreducible randomness but at the cost of modifying the
Schr\"odinger equation~\cite{GRW1986}.  More recently, 't~Hooft
has argued for a fully deterministic cellular-automaton
substrate~\cite{tHooft2016}, while the Harrigan--Spekkens classification
has sharpened the epistemic/ontic distinction for
hidden-variable models~\cite{Harrigan2010}.

We propose a resolution that requires none of these trade-offs.  Within
Recognition Science (RS)~\cite{RS_Axioms}, the cost functional~$J$ is
strictly convex, so every constrained optimization has a unique solution:
the dynamics are deterministic.  But any finite-resolution observer
projects the full ledger state through a lossy map, and the apparent
statistics of this projection reproduce quantum probabilities.
``Randomness'' is epistemic---about what an observer can know---not
ontological.

%% ============================================================
\section{The Cost Functional and Its Convexity}
\label{sec:cost}
%% ============================================================

\begin{axiom_rs}[A1 --- Normalization]
$J(1) = 0$.
\end{axiom_rs}

\begin{axiom_rs}[A2 --- Recognition Composition Law]
$J(xy) + J(x/y) = 2J(x)J(y) + 2J(x) + 2J(y)$ for all $x, y > 0$.
\end{axiom_rs}

\begin{axiom_rs}[A3 --- Calibration]
$J''_{\log}(0) = 1$, where $J_{\log}(t) := J(e^t)$.
\end{axiom_rs}

\begin{theorem}[Cost Uniqueness]
\label{thm:J}
The unique continuous solution is $J(x) = \tfrac{1}{2}(x + x^{-1}) - 1$
with $J(x) \geq 0$, $J(x) = 0 \iff x = 1$, and $J(x) = J(1/x)$.
\end{theorem}

\begin{proof}
Set $H(t) = 1 + J(e^t)$.  A2 becomes the d'Alembert functional equation
$H(s+t) + H(s-t) = 2H(s)H(t)$; A1 gives $H(0) = 1$; A3 gives $H''(0) = 1$.
By the Acz\'el classification~\cite{Aczel1966}, the continuous even solutions
with $H(0) = 1$ are $H(t) = \cosh(\alpha t)$ ($\alpha \geq 0$).  Calibration
forces $\alpha = 1$, so $J(x) = \cosh(\ln x) - 1 = \tfrac{1}{2}(x + x^{-1}) - 1$.
Non-negativity follows from AM--GM ($x + x^{-1} \geq 2$, with equality iff
$x = 1$).  Reciprocal symmetry is immediate from $x \leftrightarrow x^{-1}$.
\end{proof}

\begin{theorem}[Strict Convexity]
\label{thm:convex}
$J$ is strictly convex on $(0, \infty)$.  Indeed,
\[
  J''(x) = \frac{1}{x^3} > 0 \quad \text{for all } x > 0.
\]
\end{theorem}

\begin{proof}
$J(x) = \tfrac{1}{2}(x + x^{-1}) - 1$, so $J'(x) = \tfrac{1}{2}(1 -
x^{-2})$ and $J''(x) = x^{-3}$.  Since $x > 0$, $J''(x) > 0$ everywhere
on $(0, \infty)$.
\end{proof}

\begin{remark}
Strict convexity is a \emph{consequence} of the cost axioms, not an
independent postulate.  It is what makes the dynamics deterministic.
\end{remark}

%% ============================================================
\section{Determinism: The Unique Minimizer Theorem}
\label{sec:determinism}
%% ============================================================

\begin{definition}[Constrained Problem]
\label{def:problem}
A \emph{constrained problem} is a triple $(S, J, C)$ where $S \subseteq
\mathbb{R}_{>0}$ is a non-empty convex feasible set and $C$ encodes the
constraints.  The solution is $x^* = \arg\min_{x \in S} J(x)$.
\end{definition}

\begin{theorem}[Unique Minimizer Principle]
\label{thm:unique-min}
Let $S \subseteq \mathbb{R}_{>0}$ be a non-empty convex set.  If $x^*$
and $y^*$ are both minimizers of $J$ over $S$ (i.e.,
$J(x^*) \leq J(y)$ and $J(y^*) \leq J(z)$ for all $y, z \in S$),
then $x^* = y^*$.
\end{theorem}

\begin{proof}
Suppose $x^* \neq y^*$.  Since both are minimizers, $J(x^*) = J(y^*)$
(each is $\leq$ the other).  By convexity of $S$, the midpoint
$m = (x^* + y^*)/2 \in S$.  By strict convexity of $J$
(Theorem~\ref{thm:convex}):
\[
  J(m) < \frac{1}{2}\,J(x^*) + \frac{1}{2}\,J(y^*) = J(y^*),
\]
contradicting the minimality of $y^*$.  Therefore $x^* = y^*$.
\end{proof}

\begin{corollary}[Deterministic Dynamics]
\label{cor:determinism}
At each ledger tick, the next state is the unique $J$-minimizer over the
feasible set determined by the current state plus conservation constraints.
The dynamics are fully deterministic: given the current state, there is
exactly one possible next state.
\end{corollary}

\begin{corollary}[Zero Defect is Unique]
\label{cor:det-resolution}
There exists a unique $x > 0$ with $J(x) = 0$, namely $x = 1$.
\end{corollary}

\begin{proof}
By Theorem~\ref{thm:J}: $J(1) = 0$, and $J(x) > 0$ for all $x \neq 1$
(strict AM--GM gives $x + x^{-1} > 2$ when $x \neq 1$).
\end{proof}

%% ============================================================
\section{Apparent Randomness: Finite-Resolution Projection}
\label{sec:randomness}
%% ============================================================

\begin{definition}[Observer]
\label{def:observer}
Let the ledger contain $N \geq 1$ entries, each taking values
in~$\mathbb{R}_{>0}$, so that the full ledger state space
is~$\mathbb{R}_{>0}^N$.  An \emph{observer} is specified by:
(a)~a positive integer $m$ (the number of distinguishable outputs)
and (b)~a measurable map
$\pi : \mathbb{R}_{>0}^N \to \{1, \ldots, m\}$ from the ledger state
space to the observer's output alphabet.  We call $m$ the observer's
\emph{resolution}.
\end{definition}

\begin{definition}[Projection Fiber]
\label{def:fiber}
Given observer $\pi$ with resolution $m$, the \emph{fiber of outcome $k$}
is $F_k := \pi^{-1}(k) \subseteq \mathbb{R}_{>0}^N$.  The fibers
$\{F_1, \ldots, F_m\}$ partition the ledger state space.
\end{definition}

\begin{theorem}[Every Finite-Resolution Observer is Lossy]
\label{thm:lossy}
For every observer with finite resolution $m$, the projection map $\pi$
is many-to-one: at least one fiber $F_k$ contains uncountably many
distinct ledger states.  If $\pi$ is surjective and each fiber has
positive Lebesgue measure (i.e., the observer is physically non-degenerate),
then \emph{every} fiber contains uncountably many states.
\end{theorem}

\begin{proof}
The ledger state space $\mathbb{R}_{>0}^N$ is uncountable (for $N \geq 1$,
it has the cardinality of the continuum).  The output set $\{1, \ldots, m\}$
is finite.  Since $\mathbb{R}_{>0}^N = \bigcup_{k=1}^m F_k$ is a
partition into $m$ fibers, at least one fiber must be uncountable (an
uncountable set cannot be written as a finite union of countable sets),
so $\pi$ is many-to-one.  If additionally each fiber has positive
Lebesgue measure, then each fiber is uncountable (any measurable subset
of~$\mathbb{R}^N$ with positive measure is uncountable).
\end{proof}

\begin{remark}[Discrete dynamics in a continuous ambient space]
The RS discreteness axiom (T2) forces the physical ledger entries to
occupy discrete values (algebraic in~$\varphi$).  The ambient space
$\mathbb{R}_{>0}^N$ is used above for the cardinality argument.  Even
restricting to the discrete subset of physical states, Theorem~\ref{thm:lossy}
holds whenever the number of accessible ledger states exceeds~$m$, which is
guaranteed by the astronomical size of the ledger ($N \gg 1$).
\end{remark}

\begin{remark}[Epistemic vs.\ ontic randomness]
In the Harrigan--Spekkens classification~\cite{Harrigan2010}, RS is an
\emph{epistemic} hidden-variable model: the ontic state (the full
ledger) determines all measurement outcomes, but an observer with
finite resolution can only access an epistemic state (a coarse-grained
projection).  What separates RS from other epistemic models is that the
hidden structure is \emph{derived} from the cost functional rather than
postulated.
\end{remark}

\begin{corollary}[Every Observer Sees Apparent Randomness]
\label{cor:random}
No finite-resolution observer can distinguish all ledger states.  Multiple
distinct deterministic states produce the same observation, making the
dynamics \emph{appear} random from the observer's perspective.
\end{corollary}

\begin{corollary}[Fiber Partition]
\label{cor:fiber}
For every observer with resolution~$m$, the $m$~fibers
satisfy:
\begin{enumerate}[label=(\roman*),nosep]
  \item Every ledger state belongs to exactly one fiber (partition).
  \item Each fiber is non-empty (every output is realizable).
  \item Each fiber contains infinitely many ledger states (by
  Theorem~\ref{thm:lossy}).
\end{enumerate}
The fibers $\{F_1,\ldots,F_m\}$ provide the sample space of probability
theory without additional postulates.  What probability measure to assign to
each fiber (and in particular whether it reproduces $P = |\psi|^2$) is a
further question addressed in Remark~\ref{rem:born-status}.
\end{corollary}

\begin{proof}
(i) and (ii): immediate from the definition of a function from a
non-empty domain.  (iii): Theorem~\ref{thm:lossy}.
\end{proof}

%% ============================================================
\section{Compatibility with Bell's Theorem}
\label{sec:bell}
%% ============================================================

Bell's theorem~\cite{Bell1964} proves that no \emph{local}
hidden-variable theory can reproduce all quantum-mechanical predictions.
Does this conflict with the determinism established above?

No.  The RS ledger is a \emph{global} structure: the $J$-minimization at
each tick is performed over the \emph{entire} configuration space, not
locally.  Bell's theorem rules out local hidden variables; RS has global
(non-local) hidden structure.  The distinction is crucial:

\begin{center}
\renewcommand{\arraystretch}{1.3}
\begin{tabular}{@{}p{5cm}p{5cm}@{}}
\toprule
\textbf{Local HV (ruled out)} & \textbf{RS ledger (compatible)} \\
\midrule
Each particle carries its own hidden state & The ledger is a single global object \\
Correlations explained by shared local information & Correlations emerge from global $J$-minimization \\
Predicts Bell inequality satisfied & Bell inequality can be violated (non-local minimization) \\
\bottomrule
\end{tabular}
\end{center}

\begin{proposition}[Bell Compatibility]
\label{prop:bell}
The RS framework is compatible with Bell's theorem.
\end{proposition}

\begin{proof}
Bell's theorem~\cite{Bell1964} establishes that no \emph{local realistic}
theory can reproduce all quantum correlations.  A local realistic theory
assigns to each particle a pre-existing outcome independent of the
measurement setting at the distant wing.  RS is not such a theory: the
$J$-minimization that determines the next ledger state is performed over
the \emph{entire} configuration $(x_1,\ldots,x_N)$ simultaneously, so
the outcome assigned to a subsystem (say, entries $1$ through $k$)
depends on the state of all entries, including those far away.  This is
precisely the non-local structure that Bell's proof requires to be absent
from a local hidden variable theory.  RS provides global hidden structure,
not local hidden variables, and is therefore not excluded by Bell's
theorem.  The combination of unique global minimization (determinism) and
non-local correlations is consistent with all known Bell-inequality
violations.
\end{proof}

%% ============================================================
\section{Compatibilist Free Will}
\label{sec:freewill}
%% ============================================================

If the dynamics are deterministic, is there room for free will?  Yes.
The $J$-cost landscape is not uniformly steep.  At certain configurations,
the landscape is approximately flat: multiple futures have nearly equal
cost.  At these \emph{genuine decision points}, the agent must choose among
near-degenerate options.

\begin{definition}[Decision Point]
\label{def:decision}
A \emph{decision point} is a ledger configuration at which the feasible
set contains multiple states with near-equal $J$-cost: there exist
$x, y \in S$ with $x \neq y$ and $|J(x) - J(y)| < \varepsilon$ for some
$\varepsilon > 0$.  A decision point is \emph{genuine} when
$\varepsilon$ is of order the minimum step cost $J''(1) = 1$, so that
the cost difference between alternatives is comparable to the quantum
of defect.
\end{definition}

\begin{proposition}[Gap-45 Barrier]
\label{thm:gap45}
The 8-tick epoch ($2^3 = 8$) and the 45-fold rung structure
($45 = 3^2 \times 5$) are coprime: $\gcd(8, 45) = 1$, so they
synchronize only at multiples of $\mathrm{lcm}(8, 45) = 360$.
At all other ticks, the phase-optimal and cost-optimal successors
generically differ, providing near-degenerate decision points.
\end{proposition}

\begin{proof}
$8 = 2^3$ and $45 = 3^2 \times 5$ share no common prime factor, so
$\gcd(8, 45) = 1$ and $\mathrm{lcm}(8, 45) = 360$.  By the Chinese
Remainder Theorem, a system with 8-periodic phase constraints and
45-periodic cost constraints is simultaneously constrained only at
multiples of 360.  The conclusion that this coprimality creates
genuine decision points is a structural argument: the mathematical
coprimality is proved; the physical identification of these ticks
with near-degenerate decision points is a structural hypothesis.
\end{proof}

\begin{remark}
At these non-synchronization ticks, determining which successor has
lower total cost requires evaluating a function whose inputs depend on
the entire ledger history---a computation that no subsystem of the
ledger can perform in fewer ticks than the system itself.  This provides
a structural (not merely practical) limit on external prediction of the
agent's choice.
\end{remark}

\begin{proposition}[Compatibilist Resolution]
\label{thm:compat}
Determinism and free will coexist in the RS framework:
\begin{enumerate}[label=(\roman*)]
  \item The $J$-cost landscape is fixed and objective (determinism,
  Theorem~\ref{thm:unique-min}).
  \item At genuine decision points, the agent selects among
    near-degenerate options (free will, Definition~\ref{def:decision}).
  \item The Gap-45 barrier ensures this selection cannot be predicted
    by any external computational process in fewer ticks
    (Proposition~\ref{thm:gap45}).
\end{enumerate}
This is a structural argument combining the above results; the
computational-irreducibility claim in~(iii) is a structural hypothesis
whose full proof remains open.
\end{proposition}

\begin{remark}
This is a \emph{compatibilist} position: the laws are deterministic, but
genuine choices exist.  It differs from both hard determinism (no free will)
and libertarian free will (uncaused choices).  The choice is
\emph{underdetermined} by the cost landscape at decision points---not
because the landscape is random, but because it is flat.

Conway and Kochen's Free Will Theorem~\cite{Conway2006} establishes that
if experimenters have free will (their choices are not predetermined by
prior history), then so do elementary particles.  The RS Gap-45 barrier
provides a structural mechanism for this: at decision points, the choice
is underdetermined by the ledger history, consistent with the theorem's
requirements.
\end{remark}

%% ============================================================
\section{The Born Rule from Projection}
\label{sec:born}
%% ============================================================

In standard quantum mechanics, the Born rule ($P = |\psi|^2$) is a
postulate.  In RS, its structural basis emerges from the fiber structure
of the projection map.

The structural basis for the Born rule rests on Corollary~\ref{cor:fiber}.
For any observer with resolution $m$: the $m$ fibers partition the state
space (the sample space of probability theory); the natural measure on
each fiber provides a probability for each outcome.  A higher-resolution
observer (larger $m$) has finer fibers and can distinguish more states,
corresponding to more precise measurements.  This provides the
sample-space structure of probability without additional postulates.

\begin{remark}[Status of the Born Rule derivation]
\label{rem:born-status}
Corollary~\ref{cor:fiber} shows that any finite-resolution projection
partitions the ledger state space into fibers, providing the
\emph{sample space} of probability theory without additional postulates.
This is the structural skeleton of the Born rule.  However, the
derivation of the \emph{specific} Born weight $P = |\psi|^2$ requires
more: one must show that the RS measure on the fibers (induced by the
8-tick phase weights $e^{ik\pi/4}$, $k = 0,\ldots,7$) reproduces
exactly $|\text{amplitude}|^2$.  The full explicit derivation of the
Born probability formula from the RS fiber structure and 8-tick epoch is
\emph{not} completed in this paper; it is an identified open problem.

The comparison table's claim ``Born rule is primitive: FALSE'' is therefore
properly understood as ``the Born rule has a structural basis in the RS
fiber structure (Corollary~\ref{cor:fiber}); the complete zero-postulate
derivation of $P = |\psi|^2$ is an open problem.''
\end{remark}

%% ============================================================
\section{Laplace's Demon Dissolved}
\label{sec:laplace}
%% ============================================================

Laplace's demon---an intellect that knows the complete state and can
predict all futures---is coherent in RS.  Reality IS deterministic.  But
the demon requires access to the full ledger state, not just a projection.
For any finite-resolution observer, this access is impossible.

\begin{corollary}[Laplace's Demon Dissolved]
\label{thm:laplace}
Combining Corollary~\ref{cor:determinism} (unique minimizer) and
Theorem~\ref{thm:lossy} (lossy projection):
\begin{enumerate}[label=(\roman*)]
  \item The ledger state uniquely determines all future states
    (determinism).
  \item No finite-resolution observer can access the full ledger state
    (projection is lossy).
  \item Therefore: determinism is ontologically true but epistemically
    inaccessible to any finite observer within the ledger.
\end{enumerate}
\end{corollary}

\begin{remark}
This dissolves the Laplace demon: determinism does not imply predictability.
The demon is logically possible but physically unrealizable for any
observer within the ledger.
\end{remark}

%% ============================================================
\section{Discussion}
\label{sec:discuss}
%% ============================================================

\subsection{Summary of the Resolution}

\begin{center}
\renewcommand{\arraystretch}{1.3}
\begin{tabular}{@{}p{4.5cm}p{9cm}@{}}
\toprule
\textbf{Claim} & \textbf{Status in RS} \\
\midrule
``Reality is deterministic'' & TRUE: unique $J$-minimizer at each tick \\
``Observations are random'' & TRUE: projection through finite resolution is lossy \\
``Quantum randomness is fundamental'' & FALSE: it is epistemic (about observation), not ontological \\
``Bell's theorem rules out determinism'' & FALSE: it rules out \emph{local} determinism; RS is global \\
``Free will is an illusion'' & FALSE: genuine decision points exist at Gap-45 \\
``Born rule is primitive'' & PARTIALLY RESOLVED: the sample-space structure
(fiber partition) emerges without postulates; the full derivation of
$P = |\psi|^2$ from the fiber measure remains an open problem \\
\bottomrule
\end{tabular}
\end{center}

\subsection{Comparison with Other Interpretations}

\paragraph{Bohmian mechanics.}  Like RS, Bohmian mechanics restores
determinism with non-local hidden variables.  But Bohm introduces a
\emph{guiding equation} as an additional postulate.  RS derives everything
from the single cost functional---no guiding equation is needed.

\paragraph{Many-worlds.}  Everett's interpretation preserves determinism
by having all outcomes occur in different branches.  RS avoids ontological
proliferation: there is one ledger with one deterministic evolution.
The ``branches'' are the fibers of the projection map---mathematical
subsets, not physical worlds.

\paragraph{GRW and collapse models.}  Collapse theories introduce
irreducible randomness via stochastic modifications of the
Schr\"odinger equation.  RS requires no such modification: the apparent
randomness arises entirely from finite-resolution projection.

\paragraph{'t~Hooft's cellular automaton.}
't~Hooft~\cite{tHooft2016} also proposes a deterministic substrate
beneath quantum mechanics, arguing that a cellular automaton at the
Planck scale can reproduce quantum field theory in an appropriate
limit.  RS shares 't~Hooft's deterministic commitment but differs in
two respects: (i)~RS derives the discrete structure from the cost
functional (the automaton is not postulated but \emph{forced}), and
(ii)~apparent randomness in RS arises from finite-resolution
projection rather than from information loss at the Planck scale.

\subsection{Falsifiable Predictions}

\begin{enumerate}
  \item \textbf{No irreducible randomness.}  If an experiment ever
  demonstrates randomness that cannot be explained by projection through
  finite resolution---for example, fundamental randomness in a system
  with access to the full ledger state---the framework is refuted.

  \item \textbf{Gap-45 decision structure.}  The framework predicts that
  genuine decision points are associated with the Gap-45 periodicity.
  Neurophysiological experiments could test whether decision timing
  correlates with period-360 structure.

  \item \textbf{Born-rule deviations.}  At sufficiently high resolution
  (approaching the fundamental scale), the Born rule should show
  deviations reflecting the discrete fiber structure.
\end{enumerate}

%% ============================================================
\section{Conclusion}
\label{sec:conclusion}
%% ============================================================

The determinism--randomness debate is dissolved, not by choosing a side,
but by proving that both sides are correct about different levels of
description:
\begin{itemize}
  \item \textbf{Ontologically}, the universe is deterministic.  The strict
  convexity of the cost functional $J$ guarantees a unique minimizer at
  every ledger update.

  \item \textbf{Epistemically}, observations appear random.  Finite-resolution
  projection is provably lossy, and the resulting fiber structure
  reproduces quantum statistics.

  \item \textbf{Agentially}, free will is real.  At genuine decision points
  (flat $J$-landscape, Gap-45 barrier), the agent must choose among
  near-degenerate futures, and no external system can predict the choice.
\end{itemize}

All three levels coexist without contradiction within a single framework
derived from one cost functional and zero adjustable parameters.

%% ============================================================
\begin{thebibliography}{99}

\bibitem{RS_Axioms}
J.~Washburn, ``The Algebra of Reality: A Recognition Science Derivation
of Physical Law,'' \emph{Axioms} \textbf{15}(2), 90 (2026).
\texttt{doi:10.3390/axioms15020090}

\bibitem{Bell1964}
J.~S.~Bell, ``On the Einstein Podolsky Rosen Paradox,''
\emph{Physics} \textbf{1}, 195 (1964).

\bibitem{Bohm1952}
D.~Bohm, ``A Suggested Interpretation of the Quantum Theory in Terms of
`Hidden' Variables,'' \emph{Physical Review} \textbf{85}, 166 (1952).

\bibitem{Everett1957}
H.~Everett III, ```Relative State' Formulation of Quantum Mechanics,''
\emph{Reviews of Modern Physics} \textbf{29}, 454 (1957).

\bibitem{GRW1986}
G.~C.~Ghirardi, A.~Rimini, and T.~Weber, ``Unified dynamics for
microscopic and macroscopic systems,'' \emph{Physical Review D}
\textbf{34}, 470 (1986).

\bibitem{Laplace1814}
P.-S.~Laplace, \emph{Essai philosophique sur les probabilit\'es} (1814).

\bibitem{Conway2006}
J.~Conway and S.~Kochen, ``The Free Will Theorem,''
\emph{Foundations of Physics} \textbf{36}, 1441 (2006).

\bibitem{Aczel1966}
J.~Acz\'el, \emph{Lectures on Functional Equations and Their
Applications} (Academic Press, 1966).

\bibitem{Born1926}
M.~Born, ``Zur Quantenmechanik der Sto\ss vorg\"ange,''
\emph{Zeitschrift f\"ur Physik} \textbf{37}, 863 (1926).

\bibitem{tHooft2016}
G.~'t~Hooft, \emph{The Cellular Automaton Interpretation of Quantum
Mechanics} (Springer, 2016).

\bibitem{Harrigan2010}
N.~Harrigan and R.~W.~Spekkens, ``Einstein, Incompleteness, and the
Epistemic View of Quantum States,''
\emph{Foundations of Physics} \textbf{40}, 125 (2010).

\end{thebibliography}

\end{document}
